Large language models (LLMs) are increasingly consulted as an accessible, low-cost first point of contact for mental health concerns, particularly in regions where professional services are scarce. This dissertation investigates whether LLM responses to standardised mental health queries vary systematically with the geographic location implied by the user, in ways that disadvantage users in lower-income and lower-infrastructure settings. A dataset of 1{,}120 responses was collected from seven open-weight and commercial LLMs across 20 cities spanning four World Bank income categories and six WHO regions, using a fixed battery of eight mental health scenarios.

Rather than relying on unvalidated automated keyword coding -- a common practice in the LLM-audit literature -- this dissertation adopts a mixed-methods design in which every automated outcome measure is validated against independent human coding before being used in a confirmatory statistical model. Two of the three originally proposed outcome measures failed this validation on the first construction. In one case, an outcome intended to measure whether responses were appropriately localised to the user's stated country was found to contain a circular component that produced a result running in the opposite direction to the true pattern in the data -- a finding that is itself reported as a methodological contribution, since the flaw survived agreement between two independent automated coding rules and was caught only by comparison against human judgement.

Following remediation, three pre-registered hypotheses were tested, each with a validated outcome measure (F1 $\geq$ 0.70) and a single controlled statistical model: (H1) response actionability increases with the income category of the stated location; (H2) response localisation increases with the documented availability of a national crisis service; and (H3) the odds of a response recommending religious or spiritual support vary by WHO region. All three hypotheses were supported. H1 and H2 were additionally robust to analyses with an effective sample size of 20 cities and to the exclusion of any single model. Because the fitted H3 model adjusts for model identity and scenario but not income, its regional estimates should not be interpreted as independent of income.

A fourth, unplanned finding emerged from auditing the raw response text: several models offer specific, visibly fabricated crisis-helpline telephone numbers -- following recognisable patterns such as the North American cinematic ``555'' block or sequential digit runs -- at a rate that rises sharply as documented crisis-service availability falls, from approximately 1\% in countries with an established national line to 34\% in countries with no documented service. This is reported as a floor estimate because pattern-matching cannot detect a fabricated number that appears plausible.

The dissertation concludes that geographic and infrastructural context measurably shape the mental health guidance produced by current LLMs, that automated content-coding pipelines require independent human validation before their outputs can be trusted, and that the fabrication of specific, actionable-looking crisis contacts in underserved regions constitutes a distinct and policy-relevant safety concern warranting further investigation.
